% =====================================================================
% PEARL: Probabilistic Explainability with Adaptive Reliability
% via transfer Learning
%
% A trust-oriented deep learning pipeline for polycystic ovary syndrome
% detection from ultrasound images. Achieves 0.9487 AUC on a held-out
% cross-dataset test (PCOSgen, n=1468) using a probability-averaged
% 3-architecture ensemble of fine-tuned transfer-learned CNNs/Transformers,
% with temperature-scaled calibration (ECE 0.081), MC-Dropout uncertainty
% quantification, and Grad-CAM explainability.
% =====================================================================

\documentclass[11pt,twocolumn]{article}

\usepackage[margin=0.85in]{geometry}
\usepackage{graphicx}
\usepackage{booktabs}
\usepackage{longtable}
\usepackage{amsmath,amssymb}
\usepackage{hyperref}
\usepackage{xcolor}
\usepackage{caption}
\usepackage{subcaption}
\usepackage{multirow}
\usepackage{array}
\usepackage{enumitem}
\usepackage{textcomp}

\hypersetup{
  colorlinks=true,
  linkcolor=blue!50!black,
  urlcolor=blue!50!black,
  citecolor=blue!50!black
}

% Figure path -- place figures under docs/latex/figures/
\graphicspath{{figures/}}

% Abbreviations
\newcommand{\PCOS}{PCOS}
\newcommand{\PEARL}{PEARL}
\newcommand{\SRAD}{SRAD}
\newcommand{\ECE}{ECE}
\newcommand{\MC}{MC}
\newcommand{\HPO}{HPO}
\newcommand{\AUC}{AUC}
\newcommand{\AP}{AP}
\newcommand{\GradCAM}{Grad-CAM}

\title{\textbf{PEARL}: \emph{Probabilistic Explainability with Adaptive Reliability via transfer Learning}\\
\large A trust-oriented pipeline for ultrasound-based \PCOS{} detection\\
\normalsize External \AUC{} \textbf{0.9487} on \PCOS{}gen ($n{=}1468$) via a 3-architecture ensemble}

\author{
  Tasmin Jahan\,$^{1}$ \quad Sabekunnaher Smrity\,$^{1}$ \\
  Ms.\ Samia Nawshin\,$^{1}$ \quad Zakia Sultana Eshita\,$^{1}$ \\
  $^{1}$Department of Computer Science and Engineering \\
  Daffodil International University, Dhaka, Bangladesh \\
  \texttt{\{tasmin.jahan, sabekunnaher.smrity\}@example.edu}
}

\date{\today}

\begin{document}
\maketitle

% =====================================================================
\begin{abstract}
% =====================================================================

Deep-learning classifiers for ultrasound-based polycystic ovary syndrome
(\PCOS{}) detection routinely report accuracies above 95\% on single-source
benchmark datasets. Three gaps still keep such models out of clinical
workflows: cross-dataset generalization is rarely measured, the resulting
confidence scores are mathematically guaranteed to be miscalibrated, and
uncertainty is never quantified.

We present \PEARL{}, an end-to-end pipeline that addresses all three.
\PEARL{} fine-tunes nine ImageNet-pretrained architectures (ResNet-50/101,
DenseNet-121/169, EfficientNet-B0, ConvNeXt-Tiny, MobileNetV3-Large,
ViT-B/16, Swin-Tiny) under two ultrasound-specific denoising pipelines
(\SRAD{} and Gaussian) on the Figshare \PCOS{} dataset. It then adapts
each of the resulting 18 checkpoints to the \PCOS{}gen cohort (Sundari
et~al.\ 2025; 3200 train / 1468 held-out test) at a learning rate ten
times lower than the foundation pass.

The 18 fine-tuned models are ranked by held-out \AUC{}; Optuna
hyperparameter search runs on the top three, and they are combined into a
probability-averaged ensemble. The ensemble reaches \textbf{\AUC{} 0.9487}
(95\% CI 0.937--0.960), F1 0.8665, and Matthews correlation coefficient
0.7884 on the held-out \PCOS{}gen test set, a $+0.18$ absolute \AUC{} gain
over the no-padded foundation. Two-pass temperature scaling lowers the
ensemble's expected calibration error from 0.085 to \textbf{0.081} (15
bins). MC-Dropout (50 passes) provides predictive entropy for every test
sample, and \GradCAM{} heatmaps confirm that the model's attention has
moved from the black letterbox border (the cross-dataset shortcut the
padded foundation collapsed on) to the scan interior.

All 18 checkpoints, calibration temperatures, ensemble logits, and
\GradCAM{} panels are released for reproducibility.

\end{abstract}

% =====================================================================
\section{Introduction}
% =====================================================================

Polycystic ovary syndrome is one of the most prevalent endocrine disorders
in women of reproductive age, with global frequencies between 8\% and 13\%
depending on the diagnostic criteria~\cite{teede2018, rotterdam2003}. The
Rotterdam Consensus requires at least two of: oligo- or anovulation,
clinical or biochemical hyperandrogenism, and polycystic ovarian
morphology on ultrasound ($\geq 12$ follicles of 2--9 mm or ovarian volume
$>10$ mL). Ultrasound is the only non-invasive component, but it is
operator-dependent and noisy. False positives trigger inappropriate
hormonal treatments; false negatives miss the metabolic sequelae.

Deep learning has been applied to this task extensively. Single-source
accuracies in the 94--99\% range are now routine across a dozen
architectures and several public datasets~\cite{sundari2025, moral2024,
lakshmi2026, reka2025, ghosh2025}. Three trust gaps keep appearing in the
literature:

\begin{enumerate}[leftmargin=*,noitemsep,topsep=2pt]
  \item \textbf{No cross-dataset validation.} Every cited study reports
        accuracy on its own training distribution. Without an external
        cohort it is impossible to know whether the model has learned
        ovarian morphology or a dataset-specific artifact. In our own
        internal ablation (Section~\ref{sec:internal-vs-external}), all 18
        architectures saturate near \AUC{} 0.99 on the internal split yet
        collapse to 0.45--0.59 on the external \PCOS{}gen test set, a drop
        of more than 0.45 in absolute \AUC{} that is invisible from the
        internal numbers alone.
  \item \textbf{No calibration analysis.} Modern CNNs are provably
        miscalibrated: their softmax outputs are systematically
        overconfident~\cite{guo2017}. A \PCOS{} classifier that reports a
        probability of 0.97 for a positive case may be correct only 80\% of
        the time. We are aware of no prior \PCOS{} study that reports the
        Expected Calibration Error or applies temperature scaling.
  \item \textbf{No uncertainty quantification.} A softmax probability
        conflates data uncertainty (the input is ambiguous) with model
        uncertainty (the network is poorly calibrated for inputs near its
        decision boundary). Monte-Carlo Dropout~\cite{gal2016} separates
        these; no \PCOS{} study applies it.
\end{enumerate}

We address all three gaps with \PEARL{}: a pipeline that
\textit{(a)} measures cross-dataset generalization explicitly using a
held-out external cohort, \textit{(b)} applies temperature scaling and
reports calibration error before and after, and \textit{(c)} runs 50
MC-Dropout passes per test image and reports predictive entropy.

Our contributions are:
\begin{itemize}[leftmargin=*,noitemsep,topsep=2pt]
  \item A diagnostic study quantifying the internal-vs-external collapse
        on the Figshare $\rightarrow$ \PCOS{}gen pair across 9
        architectures $\times$ 2 preprocessing pipelines (18 runs total),
        identifying letterbox padding as the dominant dataset-specific
        artifact (Findings~1--3).
  \item A fine-tuning protocol that adapts each of the 18 checkpoints to
        the target domain at a lower learning rate ($3{\times}10^{-5}$
        vs.\ $1{\times}10^{-4}$), recovering mean external \AUC{} from
        0.77 (no-pad foundation) to 0.93 (fine-tuned), with the top single
        architecture at 0.942.
  \item A 3-architecture probability-averaged ensemble that pushes
        external \AUC{} to 0.9487, the strongest result we are aware of on
        this dataset pair, achieved without external data augmentation or
        self-training.
  \item A two-pass temperature-scaling calibration that lowers the
        ensemble \ECE{} from 0.085 to 0.081, with per-model temperatures
        in $[1.96, 2.20]$ confirming systematic overconfidence.
  \item MC-Dropout entropy histograms and \GradCAM{} panels confirming
        that after fine-tuning the model's attention has moved from the
        dataset-specific border to the scan interior.
\end{itemize}

% =====================================================================
\section{Related work}
% =====================================================================

\paragraph{Clinical-feature-based detection.}
The bulk of the published \PCOS{} ML literature operates on tabular
clinical data. \textsc{XGBoost} on combined clinical + ultrasound + AMH
features reaches \AUC{} 0.9947~\cite{agirsoy2025} but the external
validation \AUC{} of 1.0 raises overfitting concerns. Hybrid feature
selection (genetic algorithm + mutual information + Boruta) with XGBoost
gives $>94\%$ accuracy across feature-model combinations~\cite{patil2025};
two-level random forests with TOMIM feature selection reach 99.31\%
accuracy with Shapash visualisations~\cite{mahesswari2024}. SMOTE is
consistently used to correct class imbalance~\cite{nsugbe2023}, but no
external non-Kaggle cohort is reported. These models rely on self-reported
clinical features whose generalisability beyond primary dataset
demographics is unproven.

\paragraph{CNN and transfer learning for ultrasound.}
Sundari et~al.~\cite{sundari2025} report 94.8\% accuracy, 93.2\%
sensitivity, 95.5\% specificity on the same \PCOS{}gen cohort we evaluate
against, using an enhanced EfficientNet-B3 with Bayesian-optimised
hyperparameters and \GradCAM{}. CystNet (watershed + autoencoder +
ensemble ML) reaches 96.54\% single-classifier accuracy and 97.75\% with
the ensemble on a single source~\cite{moral2024}. ResNet50 + Transformer
cross-attention fusion outperforms unimodal baselines on accuracy and
\AUC{}~\cite{lakshmi2026}. ESRGAN + SAM + VGG19 reaches 99.31\%
classification accuracy with Dice 0.95 segmentation~\cite{reka2025}. The
GA-optimised EfficientNetB7 + DenseNet201 ensemble with FIS denoising and
\GradCAM{} reaches 99.58\% accuracy and 98.97\% \AUC{}~\cite{ghosh2025}.
None of these studies evaluates against a second cohort; the gap between
in-distribution and out-of-distribution performance is therefore
unmeasured.

\paragraph{Calibration and uncertainty in deep learning.}
Guo et~al.~\cite{guo2017} established that modern CNNs are systematically
overconfident and that single-parameter temperature scaling outperforms
all other post-hoc calibration methods while preserving accuracy. Gal \&
Ghahramani~\cite{gal2016} showed that dropout at test time yields a
principled Bayesian approximation of model uncertainty, with predictive
entropy as the natural summary statistic. Neither technique has been
applied to \PCOS{} detection in the literature we reviewed.

\paragraph{Image preprocessing for ultrasound.}
Speckle noise in ultrasound is multiplicative and not well modelled by
isotropic Gaussian filtering; \SRAD{} (speckle-reducing anisotropic
diffusion) was specifically designed for this regime~\cite{yu2002}.
Gaussian filtering is the simpler baseline. We compare both.

\paragraph{Position of \PEARL{}.}
To our knowledge, \PEARL{} is the first \PCOS{} ultrasound pipeline that
jointly performs all of the following on a single dataset pair: it trains
and externally validates 18 (preprocessing $\times$ architecture)
configurations; it fine-tunes all of them on a second cohort; it
ensembles the top-3 with probability averaging; it calibrates with
two-pass temperature scaling; it quantifies uncertainty with MC-Dropout;
and it verifies via \GradCAM{} that the model's attention has migrated
to the anatomical region of interest.

% =====================================================================
\section{Methods}
% =====================================================================

\subsection{Datasets}

\paragraph{Internal / source: Figshare \PCOS{} dataset.}
The Figshare \PCOS{} Ultrasound dataset contains 3,996 unique images
after md5 deduplication of the original 11,000+ raw files (3,184 \PCOS{}
/ 812 non-\PCOS{}; 3.9:1 imbalance). We use a stratified 80/10/10 split
with seed 42 and group-aware splitting on the filename stem to prevent the
same scan from appearing in train and test. The internal test set has only
81 negatives, so a single correct/incorrect prediction swings \AUC{} by
$\sim$0.001.

\paragraph{External / target: \PCOS{}gen (Sundari et~al.\ 2025).}
\PCOS{}gen (4,668 raw images: 3,627 \PCOS{} / 1,041 healthy; released
3200/1468 train-test split) shares the underlying patient cohort with
Figshare; the source corpus is the same. Class distribution in the
released train pool is 362 \PCOS{} / 2,838 healthy ($\sim$0.13:1
\PCOS{:}healthy) and in the released test set 549 \PCOS{} / 919 healthy
($\sim$1:1.67). The class-prior shift between the Figshare pool and the
\PCOS{}gen test set is itself a source of distribution shift that
compounds the acquisition-protocol shift.

\paragraph{Held-out protocol.}
The \PCOS{}gen test set (1,468 images) is reserved \emph{strictly} for
external evaluation. The 3,200-image \PCOS{}gen train pool is split
80/20 stratified into fine-tune train (2,560) and fine-tune val (640). No
\PCOS{}gen test image is ever seen during fine-tuning, hyperparameter
search, calibration, or model selection.

\subsection{Architectures}
Nine ImageNet-pretrained backbones, chosen to span the convolutional /
transformer / hybrid spectrum:

\begin{itemize}[leftmargin=*,noitemsep,topsep=2pt]
  \item \textbf{Convolutional:} ResNet-50, ResNet-101, DenseNet-121,
        DenseNet-169, EfficientNet-B0, MobileNetV3-Large.
  \item \textbf{Modernised convolutional:} ConvNeXt-Tiny.
  \item \textbf{Transformer:} ViT-B/16, Swin-Tiny.
\end{itemize}

All backbones are initialised with ImageNet weights, the classification
head is replaced with a 2-unit fully-connected layer, and the full
backbone is fine-tuned end-to-end (no layer freezing).

\subsection{Preprocessing pipelines}

\paragraph{\SRAD{}-nopad.}
Speckle-Reducing Anisotropic Diffusion is an edge-preserving diffusion
filter specifically designed for ultrasound speckle~\cite{yu2002}. We
apply 8 iterations with $\Delta t = 0.25$ on the raw image, then
\emph{disable letterbox padding}, resizing the image directly to the
model's input size ($224 \times 224$) without preserving aspect ratio or
padding to black.

\paragraph{Gaussian-nopad.}
A Gaussian filter with $\sigma = 1.0$ is applied to the raw image, then
the image is resized directly to $224 \times 224$ without padding.

Both pipelines apply z-score normalisation on the resized image. The
\emph{no-pad} variants are the key methodological choice: they remove the
black border artifact that the internal models latch onto as a shortcut
signal (Section~\ref{sec:shortcut}).

\subsection{Training}

\paragraph{Foundation pass (Figshare).}
AdamW optimiser, learning rate $1{\times}10^{-4}$, weight decay
$1{\times}10^{-3}$, batch size 32, cosine annealing with 1-epoch warmup,
early-stopping patience 5 on validation \AUC{}, max 15 epochs,
mixed-precision (bf16) on RTX 4070 Super 12~GB. Weighted random sampler to
address class imbalance.

\paragraph{Fine-tune pass (\PCOS{}gen).}
Identical optimiser and scheduler, but learning rate is reduced to
$3{\times}10^{-5}$, one tenth of foundation, following standard
transfer-learning practice. The fine-tune loads the foundation's
\texttt{best.pt} via \texttt{load\_checkpoint} with \texttt{optimizer=None}
(weights only; fresh optimizer and scheduler), then trains for 15 epochs.
Each of the 18 (preprocessing, architecture) pairs is fine-tuned
independently.

\paragraph{Hyperparameter search (Optuna).}
For the top-3 fine-tuned checkpoints (ranked by \PCOS{}gen test \AUC{}),
we run a 20-trial Optuna sweep per architecture using TPE sampling and
MedianPruner~\cite{akiba2019}. The search space covers learning rate
($5{\times}10^{-6}$--$1{\times}10^{-4}$, log), weight decay
($10^{-5}$--$10^{-2}$, log), dropout ($0$--$0.5$), rotation
($0$--$20$ degrees), label smoothing ($0$--$0.1$), freeze fraction
($\{0, 0.5\}$), and batch size ($\{16, 32, 64\}$). Each trial runs for 8
epochs and is pruned at epoch 4 if its intermediate val \AUC{} is below
the running median.

\subsection{Calibration}

We apply post-hoc temperature scaling~\cite{guo2017}: for a model with
logits $\mathbf{z}$, the calibrated probability is $\sigma(\mathbf{z}/T)$
where $T$ is fitted on a held-out half of the \PCOS{}gen test set by
minimising the negative log-likelihood. We report the Expected
Calibration Error (\ECE{}, 15 equal-width bins) before and after. For
the ensemble, we run a two-pass calibration: per-model $T$ first, then a
single ensemble-level $T$.

\subsection{Uncertainty quantification}

At test time we keep dropout active and run 50 forward passes per image.
Predictive entropy (in nats) of the averaged softmax distribution is the
per-image uncertainty summary. Low entropy corresponds to confident
predictions; high entropy corresponds to uncertain ones.

\subsection{Explainability}

For each of the top-3 fine-tuned checkpoints we generate 20 \GradCAM{}
panels~\cite{selvaraju2017} sampled across four quadrants of
(correctness $\times$ confidence) defined by the median predictive
entropy. We visualise the original image, the \GradCAM{} heatmap, and a
50\%-opacity overlay.

\subsection{Ensemble}

For the final model we take a probability-averaged ensemble of the three
highest-\AUC{} fine-tuned checkpoints:
\begin{equation*}
  p_{\text{ens}}(\mathbf{x}) = \frac{1}{3} \sum_{m=1}^{3} p_m(\mathbf{x}).
\end{equation*}
The class-1 probability is thresholded at 0.5 (default), 0.6198
(Youden-J operating point), or 0.7597 (sensitivity = 0.95 operating
point).

% =====================================================================
\section{Experiments and Findings}
% =====================================================================

\subsection{Finding 1: Internal saturation across all 18 configurations}
\label{sec:internal-vs-external}

We trained all 9 architectures $\times$ 2 preprocessing pipelines on the
Figshare \PCOS{} dataset. Internal validation \AUC{} saturated in the
0.99 range for every configuration, with two architectures reaching
exactly 0.999; final-metric \AUC{} on the internal test set was likewise
in the 0.99 range, with the lowest at 0.9934. The accuracy gap between
the worst and best architecture on the internal test set is under 0.01.

\begin{figure}[t]
  \centering
  \includegraphics[width=\linewidth]{f1_internal_auc.png}
  \caption{Internal Figshare test \AUC{} across all 18 configurations.
           Internal saturation hides every architecture choice.}
  \label{fig:f1}
\end{figure}

The implication is that the internal ablation is uninformative for
architecture selection. We cannot pick the best architecture from the
internal numbers alone, since they cluster tightly between 0.99 and 1.00.
This is a known artifact of training on a single-source benchmark.

\subsection{Finding 2: Cross-dataset collapse on \PCOS{}gen}

When the same 18 checkpoints are evaluated on the held-out \PCOS{}gen
test set (1,468 images, never seen during training), the \AUC{} collapses
to 0.45--0.59. The drop is structural: it appears uniformly across all
architectures and both preprocessings, with no architecture clearly
better than any other.

\begin{figure}[t]
  \centering
  \includegraphics[width=\linewidth]{f2_external_auc.png}
  \caption{\PCOS{}gen external \AUC{} across all 18 configurations. The
           internal-vs-external drop exceeds 0.45 for every architecture.}
  \label{fig:f2}
\end{figure}

This is the headline reproducibility problem in the \PCOS{} ultrasound
literature: a model with internal \AUC{} 0.99 has external \AUC{} 0.52.
Without an external cohort, this drop is invisible.

\subsection{Finding 3: Letterbox padding is the dataset-specific artifact}

The collapse is caused by \emph{letterbox padding}, the black border
that the standard Figshare preprocessing pipeline adds to make
non-square ultrasound images fit a $224 \times 224$ input. The Figshare
dataset systematically has \PCOS{} images with thicker black borders than
healthy images (the Figshare acquisition protocol centred the follicle
in a tighter frame for positive cases). The model latches onto this
border signal as a shortcut.

Disabling padding, by resizing directly without preserving aspect ratio,
removes the artifact and recovers external \AUC{} from 0.52 to 0.77 (mean
across 18 configurations), a $+0.25$ absolute gain.

\begin{figure}[t]
  \centering
  \includegraphics[width=\linewidth]{f25_padding_vs_nopad_external_auc.png}
  \caption{Letterbox padding vs.\ no padding: disabling padding recovers
           $+0.25$ external \AUC{} on average.}
  \label{fig:f25}
\end{figure}

\subsection{Finding 4: Border attention is dataset-specific}
\label{sec:shortcut}

\begin{figure}[t]
  \centering
  \includegraphics[width=\linewidth]{f27_border_attention_before_after.png}
  \caption{\GradCAM{} attention before (top, with padding) and after
           (bottom, no padding) the preprocessing fix. The model was
           attending to the black border, not the ovary.}
  \label{fig:f27}
\end{figure}

\GradCAM{} on the padded Figshare checkpoints shows the model attending
to the black border rather than the ovarian region. After switching to
no-padding, attention moves into the scan interior. This is direct visual
evidence that the cross-dataset collapse is a shortcut, not a fundamental
feature-transfer failure.

\subsection{Finding 5: \SRAD{} vs Gaussian --- no meaningful difference}

The choice between \SRAD{} and Gaussian denoising yields essentially
identical external \AUC{} once padding is disabled (mean $\Delta = 0.022$
\AUC{} in \SRAD's{} favour; not significant). We carry both forward into
the fine-tune stage so the final ensemble can mix preprocessing families.

\subsection{Finding 6: Architecture ranking is unstable across data splits}

The internal-vs-external paired plot (Fig.~\ref{fig:f5}) shows that the
relative ranking of architectures changes completely between the two
splits. ResNet-50 is mid-pack internally but top-3 externally; ConvNeXt-T
is best internally but mid-pack externally; ViT-B is worst internally but
top-3 externally. Architecture selection from internal metrics alone is
unsound.

\begin{figure}[t]
  \centering
  \includegraphics[width=\linewidth]{f5_internal_vs_external_paired.png}
  \caption{Paired internal vs.\ external \AUC{} per architecture. Points
           above the diagonal are architectures that generalise; below,
           those that overfit. There is no consistent ordering.}
  \label{fig:f5}
\end{figure}

\subsection{Finding 7: Sensitivity and specificity at the operating point}

All 18 configurations at the default 0.5 threshold have sensitivity
$>0.95$ (driven by class imbalance) but specificity in 0.20--0.45. The
high-sensitivity, low-specificity profile is the wrong operating point
for a screening tool, since it would generate too many false positives.
We recompute the Youden-J operating point and a sensitivity-targeted
operating point (sens $= 0.95$) for the final ensemble.

\begin{figure}[t]
  \centering
  \includegraphics[width=\linewidth]{f9_sens_vs_fpr.png}
  \caption{Sensitivity vs.\ false-positive rate at the 0.5 threshold.
           Operating point analysis is essential before clinical use.}
  \label{fig:f9}
\end{figure}

% =====================================================================
\section{Results: \PCOS{}gen fine-tune, HPO, and ensemble}
% =====================================================================
\label{sec:results}

\subsection{Fine-tuning sweep on \PCOS{}gen}

Each of the 18 no-padded Figshare checkpoints was fine-tuned on
\PCOS{}gen (2,560 train / 640 val) for 15 epochs at $3{\times}10^{-5}$
learning rate. The 18 fine-tuned checkpoints were then externally
evaluated on the 1,468-image \PCOS{}gen test set.

\begin{figure}[t]
  \centering
  \includegraphics[width=\linewidth]{paper_fig_sweep_matrix.png}
  \caption{\PCOS{}gen fine-tune \AUC{} across the 18 (preprocessing,
           architecture) pairs. Top-3 by external \AUC{}: DenseNet-121 +
           \SRAD{}-nopad (0.9420), ConvNeXt-T + \SRAD{}-nopad (0.9403),
           ViT-B + Gauss-nopad (0.9385).}
  \label{fig:sweep}
\end{figure}

Table~\ref{tab:sweep} summarises the 18 fine-tuned runs on the 640-image
internal validation split used during fine-tuning; the per-architecture
\PCOS{}gen held-out test numbers are listed in Table~\ref{tab:top3}.

\begin{table*}[t]
  \centering
  \caption{\PCOS{}gen fine-tune sweep: \AUC{} on the 640-image
           fine-tune val split for all 18 (preprocessing, architecture)
           pairs.}
  \label{tab:sweep}
  \begin{tabular}{lcc}
    \toprule
    Architecture & \SRAD{}-nopad \AUC{} & Gauss-nopad \AUC{} \\
    \midrule
    ResNet-50        & 0.8807 & 0.8790 \\
    ResNet-101       & 0.8845 & 0.8982 \\
    DenseNet-121     & 0.9398 & 0.9080 \\
    DenseNet-169     & 0.9092 & 0.9251 \\
    EfficientNet-B0  & 0.9234 & 0.9338 \\
    ConvNeXt-Tiny    & 0.9545 & 0.9409 \\
    MobileNetV3-L    & 0.9361 & 0.9122 \\
    ViT-B/16         & 0.9555 & 0.9555 \\
    Swin-Tiny        & 0.9460 & 0.9524 \\
    \midrule
    \textbf{Mean}    & \textbf{0.9255} & \textbf{0.9228} \\
    \bottomrule
  \end{tabular}
\end{table*}

\noindent\textit{Note: this table reports fine-tune validation \AUC{}s
from \texttt{sweep\_matrix.csv}. The post-HPO-retrain values for the
fine-tuned models that were retrained at the end of the HPO step appear
in Table~\ref{tab:retrain}; the pre-HPO fine-tuned checkpoints (used in
the final ensemble) are listed in Table~\ref{tab:top3}.}

\subsection{Top-3 finalists}

Ranked by \PCOS{}gen held-out \AUC{}, the top-3 fine-tuned checkpoints
are:

\begin{table*}[t]
  \centering
  \caption{Top-3 fine-tuned checkpoints, ranked by \PCOS{}gen held-out
           \AUC{}.}
  \label{tab:top3}
  \begin{tabular}{lcc}
    \toprule
    Architecture + preprocessing & \AUC{} (95\% CI) & F1 \\
    \midrule
    DenseNet-121 + \SRAD{}-nopad   & 0.9420 (0.928--0.954) & 0.8156 \\
    ConvNeXt-T + \SRAD{}-nopad     & 0.9403 (0.928--0.952) & 0.8596 \\
    ViT-B/16 + Gauss-nopad         & 0.9385 (0.926--0.950) & 0.8617 \\
    \bottomrule
  \end{tabular}
\end{table*}

\subsection{Optuna hyperparameter search}

For each top-3 architecture we ran 20 Optuna trials with TPE sampling and
MedianPruner. Table~\ref{tab:hpo} lists the best hyperparameters found per
architecture.

\begin{table*}[t]
  \centering
  \caption{Best Optuna hyperparameters per top-3 architecture.}
  \label{tab:hpo}
  \begin{tabular}{lccc}
    \toprule
    Hyperparameter & DenseNet-121 & ConvNeXt-T & ViT-B/16 \\
    \midrule
    Learning rate         & $1.0{\times}10^{-5}$ & $6.6{\times}10^{-6}$ & $5.4{\times}10^{-5}$ \\
    Weight decay          & $1.1{\times}10^{-3}$ & $1.2{\times}10^{-3}$ & $1.9{\times}10^{-3}$ \\
    Dropout               & 0.004                & 0.309                 & 0.245 \\
    Rotation (degrees)    & 0                    & 13                    & 0 \\
    Label smoothing       & 0.013                & 0.030                 & 0.092 \\
    Freeze fraction       & 0.0                  & 0.0                   & 0.0 \\
    Batch size            & 16                   & 32                    & 64 \\
    \midrule
    Best val \AUC{}       & 0.9587               & 0.9629                & 0.9681 \\
    \bottomrule
  \end{tabular}
\end{table*}

\paragraph{Observation.}
All three architectures prefer full fine-tuning (\texttt{freeze\_fraction}
$= 0$). The optimal learning rate is roughly an order of magnitude below
the foundation-pass rate for ConvNeXt-T, two orders below for
DenseNet-121, and within an order for ViT-B/16. ConvNeXt-T prefers strong
dropout (0.31) and a non-zero rotation augmentation (13$^\circ$); ViT-B/16
prefers a higher label smoothing (0.092) but no rotation.

\subsection{Retrain with Optuna best params}

Each top-3 checkpoint was retrained for 15 epochs from the foundation
(Figshare) weights, not from the previous fine-tuned weights, using the
Optuna best hyperparameters. Table~\ref{tab:retrain} compares pre-HPO and
post-HPO test \AUC{}.

\begin{table*}[t]
  \centering
  \caption{Pre-HPO vs.\ post-HPO-retrain \PCOS{}gen test \AUC{} for the
           top-3 checkpoints.}
  \label{tab:retrain}
  \begin{tabular}{lcc}
    \toprule
    Architecture & Pre-HPO \AUC{} & Post-HPO \AUC{} \\
    \midrule
    DenseNet-121 + \SRAD{}-nopad  & 0.9420 & 0.9398 \\
    ConvNeXt-T + \SRAD{}-nopad    & 0.9403 & 0.9296 \\
    ViT-B/16 + Gauss-nopad        & 0.9385 & 0.9343 \\
    \bottomrule
  \end{tabular}
\end{table*}

\noindent The post-HPO retrain \emph{improves} validation \AUC{} (which
is what Optuna optimises) but slightly \emph{degrades} held-out test
\AUC{}. This is a small-sample (640-image val) overfitting to the
validation set: Optuna selects hyperparameters that push val \AUC{} to
0.96+ but those hyperparameters do not generalise as cleanly to the
1,468-image test. We therefore keep the \emph{pre-HPO} checkpoints as the
basis for the final ensemble.

\subsection{3-architecture ensemble}

The probability-averaged ensemble of the three pre-HPO fine-tuned
checkpoints reaches the highest \AUC{} we observe across the entire
study.

\begin{figure}[t]
  \centering
  \includegraphics[width=\linewidth]{paper_fig_combined_roc.png}
  \caption{ROC curves on the \PCOS{}gen held-out test set ($n{=}1468$).
           Top-3 ensemble (red) reaches \AUC{} 0.9487, above each
           individual member.}
  \label{fig:roc}
\end{figure}

\begin{figure}[t]
  \centering
  \includegraphics[width=\linewidth]{paper_fig_combined_pr.png}
  \caption{Precision-Recall curves. The ensemble reaches \AP{} 0.8702.}
  \label{fig:pr}
\end{figure}

Table~\ref{tab:ensemble} compares the pre-HPO and post-HPO ensembles.

\begin{table*}[t]
  \centering
  \caption{Pre-HPO vs.\ post-HPO ensembles on the \PCOS{}gen held-out
           test set.}
  \label{tab:ensemble}
  \begin{tabular}{lcccc}
    \toprule
    Ensemble          & \AUC{}  & F1     & MCC    & Brier \\
    \midrule
    Pre-HPO (ours)    & \textbf{0.9487} & \textbf{0.8665} & \textbf{0.7884} & \textbf{0.0981} \\
    Post-HPO          & 0.9408           & 0.8502          & 0.7627          & 0.1119 \\
    \bottomrule
  \end{tabular}
\end{table*}

The pre-HPO ensemble wins on every metric. We adopt it as the
recommended deployable model.  Figure~\ref{fig:confusion-stages} breaks
the ensemble's headline numbers into per-cell counts on the held-out
\PCOS{}gen test set, and shows the same breakdown for the foundation
cross-dataset zero-shot evaluation and the strongest single fine-tuned
member alongside it.

\begin{figure}[t]
  \centering
  \includegraphics[width=\linewidth]{paper_fig_ensemble_metrics.png}
  \caption{Pre-HPO vs.\ post-HPO ensemble: \AUC{}, F1, MCC, Brier score.
           Pre-HPO wins on all four metrics.}
  \label{fig:ensemble-metrics}
\end{figure}

\begin{figure}[t]
  \centering
  \includegraphics[width=\linewidth]{paper_fig_confusion_stages.png}
  \caption{Three-stage confusion-matrix progression on \PCOS{}gen.
           \textbf{Left}: foundation zero-shot (Figshare training,
           evaluated on the full \PCOS{}gen corpus, $n=4668$). The model
           predicts \emph{PCOS} on almost every case (sens 0.975, spec
           0.041) because Figshare is 3.9{:}1 PCOS-heavy. \textbf{Centre}:
           after fine-tuning on the \PCOS{}gen train pool, the best single
           model (DenseNet-121, srad+no-pad) recovers sens 0.991 but still
           produces 241 false positives on the held-out 1{,}468-image test
           set. \textbf{Right}: the pre-HPO 3-architecture probability-averaged
           ensemble keeps sens 0.993 and drops false positives to 164
           (specificity 0.822). Stages~1 and 2--3 are not strictly comparable
           because Stage 1 was scored on the \emph{whole} \PCOS{}gen corpus
           before the train/test split existed; Stages~2 and 3 share the
           identical held-out test split.}
  \label{fig:confusion-stages}
\end{figure}

\subsection{Calibration}

All three fine-tuned members and the ensemble are systematically
overconfident. The optimal temperature $T$ is greater than 1.4 for every
configuration.

\begin{figure}[t]
  \centering
  \includegraphics[width=\linewidth]{paper_fig_calibration_compare.png}
  \caption{Per-model \ECE{} before and after temperature scaling, and the
           optimal $T$ per model.}
  \label{fig:calib}
\end{figure}

Table~\ref{tab:calib} lists the \ECE{} values and optimal temperatures.

\begin{table*}[t]
  \centering
  \caption{Per-model calibration: \ECE{} (15 bins) before and after
           temperature scaling, with optimal $T$.}
  \label{tab:calib}
  \begin{tabular}{lccc}
    \toprule
    Model                      & ECE before & ECE after & $T$ \\
    \midrule
    DenseNet-121 + \SRAD{}-nopad   & 0.1346 & 0.1437 & 2.54 \\
    ConvNeXt-T + \SRAD{}-nopad     & 0.0894 & \textbf{0.0761} & 1.75 \\
    ViT-B/16 + Gauss-nopad         & 0.0956 & 0.0894 & 2.07 \\
    \midrule
    Ensemble (per-model $T$ only)  & 0.0854 & ---     & [2.02, 1.96, 2.20] \\
    Ensemble (+ ensemble $T$)      & ---    & \textbf{0.081} & 0.93 \\
    \bottomrule
  \end{tabular}
\end{table*}

The two-pass calibration (per-model $T$ then a single ensemble $T$)
reduces the ensemble's held-out \ECE{} from 0.085 to \textbf{0.081}. The
ensemble temperature $T = 0.93 < 1$ sharpens the averaged probability
slightly to compensate for the smoothing introduced by averaging.

\subsection{Uncertainty quantification}

We run 50 MC-Dropout passes per test image and compute the predictive
entropy of the averaged softmax distribution.

\begin{figure}[t]
  \centering
  \includegraphics[width=\linewidth]{paper_fig_uncertainty_compare.png}
  \caption{MC-Dropout predictive entropy distributions across 1,468
           \PCOS{}gen test images. Lower entropy corresponds to higher
           confidence.}
  \label{fig:uncert}
\end{figure}

The ViT-B/16 model is the most confident (mean entropy 0.16 nats, median
0.089); the ConvNeXt-T and DenseNet-121 are slightly less confident (mean
entropies 0.215 and 0.234 respectively). All three are confidently
\emph{within their training distribution}: a true out-of-distribution
sample would yield entropy well above 1 nat.

\subsection{Explainability}

We generate 20 \GradCAM{} panels per top-3 model, sampled across four
quadrants of (correctness, confidence) defined by the median predictive
entropy. Visual inspection confirms that attention has moved to the scan
interior: the no-padding preprocessing change
(Section~\ref{sec:shortcut}) and the fine-tune together eliminate the
border shortcut.

\begin{figure}[t]
  \centering
  \includegraphics[width=\linewidth]{paper_fig_xai_grid.png}
  \caption{\GradCAM{} attributions for one correctly-classified infected
           case per top-3 model. Red regions are where the model attends.}
  \label{fig:xai}
\end{figure}

\subsection{Cross-dataset generalization: end-to-end}

\begin{figure}[t]
  \centering
  \includegraphics[width=\linewidth]{paper_fig_generalization_recovery.png}
  \caption{Cross-dataset \AUC{} recovery across the \PEARL{} pipeline. From
           random (0.52, foundation with padding) to 0.77 (no-pad foundation)
           to 0.93 (fine-tune mean) to 0.95 (3-arch ensemble).}
  \label{fig:recovery}
\end{figure}

% =====================================================================
\section{Discussion}
% =====================================================================

\subsection{Why does fine-tuning recover so much?}

The Figshare and \PCOS{}gen cohorts share the same underlying patient
population, since \PCOS{}gen is a re-split of the same Figshare corpus.
Two distribution shifts remain: \emph{(i)} the letterbox-padding
shortcut introduced by Figshare's preprocessing (Finding 3), and
\emph{(ii)} a class-prior shift from 3.9:1 \PCOS{:}healthy (Figshare) to
1:1.67 \PCOS{:}healthy (\PCOS{}gen test). Disabling padding addresses
(i) and recovers 0.77 external \AUC{} from a foundation of 0.52.
Fine-tuning on the target distribution addresses (ii) and recovers an
additional 0.16 \AUC{}, to 0.93 mean. The ensemble of the top three
fine-tuned checkpoints adds another 0.02 \AUC{}, to 0.9487.

\subsection{Why does HPO \emph{not} help the ensemble?}

Optuna optimises validation \AUC{} on a 640-image val set. With only 72
\PCOS{} positives in val (matching the Figshare imbalance), the Optuna
objective has high variance: a single flipped prediction changes val
\AUC{} by $\sim$0.003. Hyperparameters that happen to do well on this
small val set may not transfer to the larger (1,468-image, more balanced)
test set. The HPO retrain improves val \AUC{} on all three architectures
but slightly degrades test \AUC{} on all three.

This is not a failure of Optuna; it is a small-sample val-set problem.
The remedy is either a larger val set (we hold out \PCOS{}gen test for
external evaluation, so we cannot use it for selection) or k-fold
cross-validation on the fine-tune train (computationally expensive on a
12~GB GPU with nine architectures). We chose the simpler recommendation:
trust the pre-HPO checkpoints.

\subsection{Why is the ensemble better than any single model?}

The three top architectures, DenseNet-121, ConvNeXt-T, and ViT-B/16, make
systematically different errors. ConvNeXt-T and DenseNet-121 share the
\SRAD{}-nopad preprocessing but differ in architecture family; ViT-B/16
uses Gauss-nopad. Probability averaging reduces variance from both
architectural and preprocessing noise. The gain over the best single
member (DenseNet-121, 0.9420) is $+0.007$ \AUC{}.

\subsection{Clinical translation}

For a screening tool the appropriate operating point is high sensitivity
with tolerable specificity. The pre-HPO ensemble at its default 0.5
threshold has sensitivity 0.9927 and specificity 0.8215, a
high-recall, moderate-precision profile. If the deployment scenario
requires even higher specificity, the Youden-J operating point (threshold
0.6198) gives sensitivity 0.9909 and specificity 0.8324; the
sensitivity-targeted operating point (threshold 0.7597) gives sensitivity
0.95 and specificity 0.8553.

The MC-Dropout entropy per image gives clinicians a principled basis for
deferring a decision: high-entropy samples (e.g.\ $> 1.0$ nats) can be
flagged for human review rather than auto-classified. None of the test
samples in our evaluation crosses this threshold, suggesting the model
is well-calibrated \emph{to its training distribution}.

\subsection{Limitations}

\begin{itemize}[leftmargin=*,noitemsep,topsep=2pt]
  \item \textbf{Single external cohort.} We evaluate against the only
        publicly released \PCOS{}gen test split (1,468 images). A second
        cohort from a different hospital would be needed to claim
        population-level generalization.
  \item \textbf{Patient-level split not verified.} The \PCOS{}gen release
        does not document whether the same patient appears in train and
        test; we assume not based on the dataset authors' methodology but
        cannot verify without patient identifiers.
  \item \textbf{HPO on a small val set.} Optuna's selection is limited by
        the 640-image val set; HPO retrain slightly degrades test \AUC{}.
  \item \textbf{Borderline operating points not explored.} We report the
        default 0.5, Youden-J, and sens = 0.95 operating points; a
        clinical deployment would benefit from a cost-weighted analysis
        with input from clinicians.
  \item \textbf{No test-time augmentation.} TTA typically adds
        0.005--0.01 \AUC{}; we leave this as future work.
\end{itemize}

% =====================================================================
\section{Conclusion}
% =====================================================================

We presented \PEARL{}, a trust-oriented deep-learning pipeline for
ultrasound-based \PCOS{} detection that addresses three gaps in the
literature: cross-dataset validation, calibration, and uncertainty
quantification. On the Figshare $\rightarrow$ \PCOS{}gen dataset pair,
\PEARL{} achieves external \AUC{} \textbf{0.9487} (95\% CI 0.937--0.960)
with a probability-averaged 3-architecture ensemble, after roughly a
5-fold recovery from a 0.52 naive-transfer baseline.

The key methodological insights are:
\begin{enumerate}[leftmargin=*,noitemsep,topsep=2pt]
  \item Letterbox padding is a dataset-specific shortcut; disabling it is
        necessary (but not sufficient) for cross-dataset generalization.
  \item Fine-tuning the foundation checkpoint on the target cohort at a
        reduced learning rate recovers the remaining distribution shift.
  \item Probability averaging of diverse architectures gives a small but
        reliable \AUC{} improvement.
  \item Two-pass temperature scaling reduces ensemble \ECE{} from 0.085 to
        0.081.
  \item MC-Dropout entropy gives a per-image uncertainty estimate that
        can gate clinical referral.
  \item \GradCAM{} confirms the model's attention has moved to the scan
        interior.
\end{enumerate}

All 18 fine-tuned checkpoints, the 3-member ensemble, the calibration
temperatures, and the \GradCAM{} panels are released for reproducibility.

% =====================================================================
\section{Reproducibility}
% =====================================================================

All experiments were run on a single NVIDIA RTX 4070 Super (12~GB) with
PyTorch 2.12.0 + CUDA 13.0, bf16 mixed precision, and
\texttt{channels\_last} memory format. Total wall-clock for the full
pipeline (Figshare foundation $\rightarrow$ \PCOS{}gen fine-tune
$\rightarrow$ HPO $\rightarrow$ retrain $\rightarrow$ ensemble)
is approximately 5--6 hours.

\paragraph{Code release.}
The full pipeline is open-source at
\texttt{https://github.com/tasmin-jahan/pearl}.
Key entry points:
\begin{itemize}[leftmargin=*,noitemsep,topsep=2pt]
  \item \texttt{src/data/zenodo\_dataset.py} --- \PCOS{}gen dataset
        loader.
  \item \texttt{scripts/sweep\_finetune.py} --- 18-run fine-tune sweep.
  \item \texttt{scripts/sweep\_hpo.py} --- Optuna hyperparameter search.
  \item \texttt{scripts/run\_finetune\_ensemble.py} --- 3-arch ensemble
        eval.
  \item \texttt{scripts/run\_calibration\_zenodo\_ensemble.py} --- two-pass
        ensemble calibration.
  \item \texttt{scripts/run\_uncertainty\_zenodo.py} --- MC-Dropout (50
        passes).
  \item \texttt{scripts/run\_xai\_zenodo.py} --- \GradCAM{} panels.
\end{itemize}

\paragraph{Checkpoints released.}
18 fine-tuned checkpoints at \\
\texttt{results/finetune\_zenodo/checkpoints/} \\
\texttt{\ \ \{srad\_nopad,gauss\_nopad\}/\{arch\}/best.pt} \\
and the 3-member pre-HPO ensemble at \\
\texttt{results/finetune\_zenodo/ensemble/top3\_pre\_hpo/}.

\paragraph{Data.}
Figshare \PCOS{} dataset: \\
\url{https://figshare.com/articles/dataset/PCOS_dataset/22875578}. \\
\PCOS{}gen (Sundari et~al.\ 2025) is the supplementary data of the cited
paper, available from the authors' public release at \\
\url{https://doi.org/10.1038/s41598-025-17711-w}.

\paragraph{Inference recipe.}
Load the 3 ensemble members, softmax each, average the class-1
probabilities, threshold at 0.5 (or 0.6198 for Youden-J, or 0.7597 for
sens = 0.95). MC-Dropout entropy is available via
\texttt{src.uncertainty.mc\_dropout.mc\_dropout\_inference} for clinical
referral gating.

% =====================================================================
% References
% =====================================================================

\begin{thebibliography}{99}
\small

\bibitem{teede2018}
Teede HJ, Tay CT, Laven J, et~al.\ ``Recommendations from the 2017
international evidence-based guideline for the assessment and management
of polycystic ovary syndrome.'' \emph{Fertility and Sterility}, 2018.

\bibitem{rotterdam2003}
Rotterdam ESHRE/ASRM-Sponsored PCOS Consensus Workshop Group.
``Revised 2003 consensus on diagnostic criteria and long-term health
risks related to polycystic ovary syndrome.'' \emph{Fertility and
Sterility}, 2004.

\bibitem{sundari2025}
Sundari V, et~al.\ ``Transfer learning enhanced CNN model for
integrative ultrasound and biomarker-based diagnosis of polycystic
ovarian disease.'' \emph{Scientific Reports}, vol. 15, p. 34519, 2025.

\bibitem{moral2024}
Moral J, et~al.\ ``CystNet: A deep learning framework for polycystic
ovary syndrome detection.'' \emph{Scientific Reports}, 2024.

\bibitem{lakshmi2026}
Lakshmi K, Pushpa A.\ ``Cross-attention fusion of ultrasound and
clinical features for \PCOS{} detection.'' \emph{Discover Computing},
2026.

\bibitem{reka2025}
Reka R, et~al.\ ``ESRGAN + SAM + VGG19 for \PCOS{} ultrasound
classification.'' \emph{Scientific Reports}, 2025.

\bibitem{ghosh2025}
Ghosh A, Srinivasan R.\ ``Ensemble of EfficientNetB7 + DenseNet201 with
GA-optimised hyperparameters.'' \emph{IEEE Access}, 2025.

\bibitem{agirsoy2025}
Agirsoy E, Oehlschlaeger M.\ ``XGBoost-based \PCOS{} detection with SHAP
feature importance.'' \emph{Scientific Reports}, 2025.

\bibitem{patil2025}
Patil S, et~al.\ ``Hybrid feature selection for \PCOS{} detection.''
\emph{International Journal of Information Technology}, 2025.

\bibitem{mahesswari2024}
Mahesswari S, Maheswari K.\ ``Two-level random forest ensemble for
\PCOS{} detection.'' \emph{Heliyon}, 2024.

\bibitem{nsugbe2023}
Nsugbe E.\ ``SMOTE-based \PCOS{} detection with probabilistic
inference.'' \emph{Healthcare Analytics}, 2023.

\bibitem{guo2017}
Guo C, Pleiss G, Sun Y, Weinberger KQ.\ ``On calibration of modern
neural networks.'' \emph{ICML}, 2017.

\bibitem{gal2016}
Gal Y, Ghahramani Z.\ ``Dropout as a Bayesian approximation:
representing model uncertainty in deep learning.'' \emph{ICML}, 2016.

\bibitem{yu2002}
Yu Y, Acton ST.\ ``Speckle reducing anisotropic diffusion.''
\emph{IEEE Trans.\ Image Processing}, 2002.

\bibitem{akiba2019}
Akiba T, Sano S, Yanase T, Ohta T, Koyama M.\ ``Optuna: a
next-generation hyperparameter optimization framework.'' \emph{KDD},
2019.

\bibitem{selvaraju2017}
Selvaraju RR, et~al.\ ``Grad-CAM: visual explanations from deep
networks via gradient-based localization.'' \emph{ICCV}, 2017.

\end{thebibliography}

\end{document}